% ============================================================
\subsubsection{Diseño del sistema}
% ============================================================

\paragraph{Modelo de casos de uso}

La Figura~\ref{fig:cu_general} presenta el diagrama general de casos de uso,
derivado de los requisitos funcionales. La Tabla~\ref{tab:traza_cu} muestra la
correspondencia entre casos de uso y requisitos.

\figura{casos_uso_general.png}{Diagrama general de casos de uso}{fig:cu_general}

\begin{table}[H]
\centering\small
\caption{Trazabilidad entre casos de uso y requisitos funcionales}\label{tab:traza_cu}
\begin{tabular}{|l|l|l|}
\hline
\textbf{Código} & \textbf{Caso de uso} & \textbf{RF} \\ \hline
CU01 & Iniciar sesión & RNF02 \\ \hline
CU02 & Gestionar paciente & RF01 \\ \hline
CU03 & Registrar historial clínico & RF02 \\ \hline
CU04 & Registrar medidas antropométricas & RF03, RF04 \\ \hline
CU05 & Registrar datos biométricos & RF05 \\ \hline
CU06 & Registrar hábitos dietéticos & RF06 \\ \hline
CU07 & Calcular GET y distribución de macronutrientes & RF07, RF08 \\ \hline
CU08 & Consultar asistente IA por paciente & RF09 \\ \hline
CU09 & Consultar asistente IA global & RF10 \\ \hline
CU10 & Cargar documento PDF & RF11 \\ \hline
CU11 & Gestionar catálogo de alimentos & RF12 \\ \hline
\end{tabular}
\end{table}
\pendiente{ajusta los códigos y nombres a los de tu diagrama}

Se describen a continuación los casos de uso centrales del prototipo, con su
flujo normal y sus flujos alternos. \pendiente{revisa que los mensajes coincidan
con los que muestra tu interfaz}

% ── Plantilla de ficha de caso de uso ───────────────────────
\newcommand{\fichaCU}[6]{%
\begin{longtable}{|p{3cm}|p{11cm}|}
\caption{Descripción del caso de uso #1}\\ \hline
\textbf{Caso de uso} & #1 \\ \hline
\textbf{Código} & #2 \\ \hline
\textbf{Actor} & Nutricionista \\ \hline
\textbf{Descripción} & #3 \\ \hline
\textbf{Precondiciones} & #4 \\ \hline
\multicolumn{2}{|c|}{\textbf{CURSO NORMAL DE EVENTOS}} \\ \hline
\multicolumn{2}{|p{14.5cm}|}{#5} \\ \hline
\multicolumn{2}{|c|}{\textbf{CURSO ALTERNO DE EVENTOS}} \\ \hline
\multicolumn{2}{|p{14.5cm}|}{#6} \\ \hline
\end{longtable}}

{\small
\fichaCU{Registrar medidas antropométricas}{CU04}
{El nutricionista registra las medidas antropométricas de un paciente y el sistema calcula el IMC y la relación cintura-cadera.}
{Sesión iniciada. Paciente seleccionado.}
{\begin{enumerate}[nosep]
 \item El nutricionista abre la sección de medidas antropométricas del paciente.
 \item Ingresa peso, talla, circunferencias de cintura y cadera, \% de grasa y masa muscular.
 \item El sistema calcula el IMC y la relación cintura-cadera y muestra su clasificación.
 \item El nutricionista guarda el registro.
 \item El sistema almacena la medición con su fecha y actualiza el historial del paciente.
\end{enumerate}}
{\textbf{A. Campos obligatorios vacíos o no numéricos:} el sistema muestra un aviso por campo y no guarda.\\
 \textbf{B. Error del servidor:} el sistema muestra un mensaje de error y conserva los datos del formulario.}

\fichaCU{Consultar asistente IA por paciente}{CU08}
{El nutricionista conversa con el asistente usando como contexto los datos clínicos del paciente seleccionado y el catálogo SMAE.}
{Sesión iniciada. Paciente seleccionado.}
{\begin{enumerate}[nosep]
 \item El nutricionista escribe una consulta en el chat del paciente.
 \item El sistema arma el contexto con los datos clínicos, antropométricos y biométricos del paciente.
 \item Si la consulta trata sobre planes, alimentos o equivalentes, el sistema añade al contexto una muestra del catálogo SMAE por grupo.
 \item El sistema envía al modelo de lenguaje el prompt del sistema, el contexto, el historial reciente y la consulta.
 \item El sistema muestra la respuesta en formato de tablas y guarda ambos mensajes en el historial.
\end{enumerate}}
{\textbf{A. Mensaje vacío:} el sistema responde «El mensaje es requerido».\\
 \textbf{B. Consulta fuera del dominio nutricional:} el asistente responde con el mensaje de rechazo definido en el prompt del sistema.\\
 \textbf{C. Límite de peticiones de la API:} el sistema informa que el servicio está saturado y pide reintentar.}

\fichaCU{Cargar documento PDF}{CU10}
{El nutricionista adjunta un documento PDF para que el asistente responda con base en su contenido.}
{Sesión iniciada.}
{\begin{enumerate}[nosep]
 \item El nutricionista adjunta un archivo PDF o de texto y escribe una consulta.
 \item El sistema valida el tipo y el tamaño del archivo (máximo 10\,MB).
 \item El sistema extrae el texto del documento.
 \item El sistema envía al modelo el texto extraído junto con la consulta.
 \item El sistema muestra la respuesta basada en el documento.
\end{enumerate}}
{\textbf{A. Tipo de archivo no permitido:} el sistema rechaza la carga.\\
 \textbf{B. Archivo mayor a 10\,MB:} el sistema rechaza la carga.\\
 \textbf{C. PDF sin texto extraíble (escaneado):} \pendiente{describe el comportamiento observado}.}
}

\paragraph{Modelo de dominio y comportamiento}

El modelo conceptual se refinó en el diagrama de clases de la
Figura~\ref{fig:clases}, con atributos, tipos y cardinalidades. Los diagramas
de secuencia de las Figuras~\ref{fig:seq_chat} y~\ref{fig:seq_pdf} describen los
dos flujos principales del asistente.

\figura{diagrama_clases.png}{Diagrama de clases de NutriAPP}{fig:clases}
\figura{secuencia_chat_paciente.png}{Diagrama de secuencia «Consultar asistente IA por paciente»}{fig:seq_chat}
\figura{secuencia_carga_pdf.png}{Diagrama de secuencia «Cargar documento PDF»}{fig:seq_pdf}
